\subsubsection{Topological sort (Kahn + DFS)}

\paragraph{Idea intuitiva.}
Un \textbf{orden topologico} de un DAG es una permutacion de los vertices tal que toda arista $u \to v$ tiene $u$ antes que $v$. Dos algoritmos clasicos:
\begin{itemize}\setlength\itemsep{0pt}
	\item \textbf{Kahn}: BFS sobre nodos con \texttt{indeg == 0}; al ``consumir'' un nodo, decrementar el indeg de sus vecinos; los que lleguen a 0 se encolan.
	\item \textbf{DFS}: hacer DFS; al terminar un nodo (estado = 2), aniadirlo al final; el orden es el reverso.
\end{itemize}
La demostracion: en un DAG siempre hay al menos un nodo con indeg 0; quitarlo mantiene el DAG; por induccion, se obtiene el orden.

\paragraph{Propiedades clave (invariantes).}
\begin{itemize}\setlength\itemsep{0pt}
	\item Existe sii el grafo es un DAG.
	\item Si al final \texttt{order.size() < V}, hay un ciclo.
	\item Kahn da el orden ``natural'' (los que pueden ir primero); DFS da el orden ``al reves'' (los que pueden ir al final primero).
	\item La cantidad de ordenes topologicos puede ser exponencial.
\end{itemize}

\paragraph{Codigo clave.}
Kahn con priority queue (orden lexicografico):
\begin{verbatim}
priority_queue<int, vector<int>, greater<int>> pq;
for (int v = 0; v < n; ++v) if (indeg[v] == 0) pq.push(v);
while (!pq.empty()) {
    int v = pq.top(); pq.pop();
    order.push_back(v);
    for (int u : adj[v]) {
        if (--indeg[u] == 0) pq.push(u);
    }
}
\end{verbatim}

\paragraph{Complejidad.}
\begin{itemize}\setlength\itemsep{0pt}
	\item $O(V + E)$ ambos.
	\item Con priority queue, $O((V + E) \log V)$.
\end{itemize}

\paragraph{Donde tocar para modificar.}
\begin{itemize}\setlength\itemsep{0pt}
	\item \textbf{Detectar ciclos}: comparar \texttt{order.size()} con \texttt{V}.
	\item \textbf{Orden lexicografico minimo}: usar priority queue en Kahn.
	\item \textbf{Topological sort con pesos}: combinar con shortest path en DAG (un solo pase en orden topologico).
	\item \textbf{Encontrar todos los ordenes}: DP sobre subsets (solo para $V$ pequeno).
\end{itemize}

\paragraph{Trampas y errores frecuentes.}
\begin{itemize}\setlength\itemsep{0pt}
	\item En el snippet, Kahn retorna orden vacio si hay ciclo (util para detectar).
	\item En DFS, no olvidar \texttt{reverse(order)}.
	\item Confundir indeg con outdeg; el orden Kahn quita primero los nodos ``fuente''.
	\item Para grafos no DAG, el algoritmo termina con \texttt{order.size() < V}.
\end{itemize}

\paragraph{Codigo.}
Ver \texttt{cpp.tex}, seccion \textit{Topological sort (Kahn + DFS)}.

\vspace{0.5em}
